\section{Threats to Validity}
\label{sec:threats}

Threats to validity are discussed below following established guidelines~\cite{wohlin2012experimentation}. Table~\ref{tab:threats-summary} summarises the principal threats and the corresponding mitigations.

\begin{table}[ht]
\centering
\caption{Summary of validity threats and mitigations.}
\label{tab:threats-summary}
\begin{tabular}{>{\raggedright\arraybackslash}p{3.5cm}>{\raggedright\arraybackslash}p{8.5cm}}
\toprule
\textbf{Threat} & \textbf{Mitigation} \\
\midrule
LLM sampling non-determinism & Five independent runs per method per phase, statistical aggregation, low sampling temperature (0.3) \\
Prompt confounding & Minimal prompts that differ only in the inclusion of specifications or source code; P4 source-code condition as upper-bound control \\
Subject-selection bias & Diverse categories across eleven classes of Apache Commons Lang; all methods of selected classes included \\
Small sample / single corpus & 312 methods (paired design); follow-up planned on service-boundary code \\
LLM familiarity with training data & P4 (source) included as a control; observed P1$\rightarrow$P3 gain inconsistent with simple training-set recall \\
Single LLM family & Replication package parameterises LLM client; cross-model replication identified as future work \\
Mutation-testing limitations & PIT default operators, equivalent-mutant filtering, complementary test-count metric reported \\
Heuristic soundness & OpenJML ESC discharge of every inferred clause (Table~\ref{tab:discharge}); 576-test verification regression suite \\
Code-derived ground truth (e.g.\ sort stability for duplicates is an implementation choice, not an abstract guarantee) & Documented as a structural limit (Section~\ref{sec:discussion}, ``Specifications underdetermined by code''); cross-implementation comparison and documentation-based recovery identified as follow-up directions \\
\bottomrule
\end{tabular}
\end{table}

\paragraph{Internal validity.} Sampling non-determinism from the LLM was mitigated through five independent runs combined with statistical analysis. Prompt confounding was addressed by keeping prompts minimal---differing only in the inclusion of specifications or source code---and by including P4 as a control. The spec-quality oracle is OpenJML's Extended Static Checker rather than human judgement, eliminating rater subjectivity at the cost of a verifier-implementation-consistency interpretation of the resulting numbers rather than a documented-intent interpretation.

\paragraph{External validity.} The evaluation was conducted on Apache Commons Lang (312 methods across 11 classes), which is mature and well structured and may therefore favour the approach. Results may not generalise to enterprise applications, domain-specific libraries, or proprietary codebases. The subset also lacks methods in the \emph{other} category (event handlers, callbacks, I/O); the strength-distribution result for that category therefore depends on broader evidence in the literature~\cite{dragan2006reverse} rather than on within-corpus data. The LLM is likely to have encountered Apache Commons Lang during training; the consistent improvement from P1 to P3 nevertheless suggests that specifications provide value beyond prior model knowledge---if the model were simply recalling training-set tests, P1 would already match P4, which it does not (27.5\% vs.\ 94.8\% mutation). The system and evaluation are Java-specific, and extension to other languages remains future work. A single LLM (Gemini~2.0) was used; the choice was made for budget reasons---the experiment requires 6{,}240 generations and was repeated several times during pilot tuning, totalling approximately 25{,}000 LLM invocations---and on the basis that Gemini~2.0 is widely available and representative of current-generation instruction-tuned LLMs. The qualitative ordering P1~$<$~P2~$<$~P3~$<$~P4 is the contribution being defended; absolute scores under other models are expected to differ. Recent multi-model evaluations of LLM test generation~\cite{schafer2023adaptive, yuan2024chattester, molinelli2025oracles} report consistent qualitative findings across GPT, Claude, and Gemini families, which provides indirect support that the ordering observed here is not Gemini-specific, but a full multi-model evaluation at the 312-method scale is left to a follow-up study and is explicitly identified as a limitation of the present article. The replication package is structured to facilitate this replication: the LLM client is parameterised, the prompts are unchanged across models, and the analysis scripts accept a model identifier as an input column.

\paragraph{Heuristic soundness.} Because inference is heuristic, individual clauses could in principle be syntactically malformed or semantically false. This threat is mitigated by discharging every inferred clause through OpenJML's Extended Static Checker (Section~\ref{subsec:validation}); a continuously executed verification suite of 576 end-to-end tests gates engine changes against this oracle, and the discharge layer is extended with \texttt{define-fun-rec} support for the \texttt{\textbackslash sum}, \texttt{\textbackslash product}, and \texttt{\textbackslash num\_of} quantifiers so that loop accumulator postconditions can be discharged rather than recorded as unsupported. Clauses that fail formal verification are flagged rather than silently retained, and the 7.1\% non-discharge residue is characterised in Section~\ref{sec:discussion} as three categorically identifiable SMT/JML tractability limits (recursive multiplicative growth, quantified invariants over mutable arrays, OpenJML for-each translation), not a failure mode of the inferrer. The verifier oracle does not measure completeness against Javadoc-documented intent; the corresponding limitation---that an inferred contract may be consistent with the implementation while still being a weaker oracle than a perfectly documented one---is the trade-off accepted in exchange for the eliminated rater-subjectivity threat.

\paragraph{Construct validity.} Mutation testing has known limitations~\cite{papadakis2019mutation}, principally equivalent mutants and the representativeness of mutation operators. The PIT default operator set~\cite{coles2016pitest} was used, and trivially equivalent mutants were excluded. Test count is reported as a complementary metric to mutation score. A second construct issue is that the spec-quality oracle is the implementation: the OpenJML discharge layer measures consistency between the inferred clause and the code that was inferred from, not consistency with an abstract API contract. For methods whose abstract contract is strictly weaker than their implementation---sort stability for duplicate keys is the canonical case---the inferred specification will be sound against the implementation but over-specified against the contract. This limit is intrinsic to code-derived inference and is discussed under ``Specifications underdetermined by code'' in Section~\ref{sec:discussion}.

\paragraph{Conclusion validity.} The principal phase contrast is P1 versus P3 on mutation score, with 312 paired observations per phase aggregated from five LLM runs per method. With paired observations and a within-pair standard deviation of 11.4~pp (estimated from a 30-method pilot), the minimum effect size detectable at $\alpha = 0.0021$ (the Bonferroni-corrected per-test threshold; see Section~\ref{sec:methodology}) and 80\% power is 2.2~pp. The observed effect is 40.7~pp, giving a post-hoc power exceeding 0.999. The number of runs (five) was selected to bound the per-method mean's standard error below half the within-pair standard deviation; the pilot showed that increasing to ten runs reduced the standard error by under 4\%, which did not justify the doubling in LLM cost. Bonferroni correction was applied across the family of 24 phase $\times$ metric contrasts, and all reported significant differences remain significant after correction.
